\uuid{3igx}
\titre{Comparer deux capteurs avec des ReLU}
\niveau{L3}
\module{OPTRN}
\chapitre{Réseaux de neurones}
\sousChapitre{Fonctions réalisées par un réseau}
\theme{ReLU, valeur absolue}
\auteur{Maxime Nguyen}
\datecreate{2026-09-25}
\organisation{AMSCC}
\difficulte{2}

\contenu{
\texte{Deux capteurs fournissent des mesures réelles $x$ et $y$. On veut calculer le score de désaccord $S(x,y)=|x-y|$. On note $\operatorname{ReLU}(t)=\max(0,t)$.}

\begin{enumerate}
\item \question{En distinguant les cas $x\geq y$ et $x<y$, démontrer l'identité
\[
|x-y|=\operatorname{ReLU}(x-y)+\operatorname{ReLU}(y-x).
\]}
\indication{Dans chaque cas, l'un des deux termes de droite est nul.}
\reponse{Si $x\geq y$, alors $\operatorname{ReLU}(x-y)=x-y$ et $\operatorname{ReLU}(y-x)=0$ ; leur somme vaut $|x-y|$. Si $x<y$, les rôles sont échangés et la somme vaut $y-x=|x-y|$.}

\item \question{Dessiner un réseau à deux neurones cachés calculant $S(x,y)$. Préciser les poids et les biais des deux neurones cachés ainsi que ceux du neurone de sortie, dont l'activation est l'identité.}
\indication{Chaque neurone caché calcule l'un des deux termes de l'identité précédente.}
\reponse{La couche cachée calcule
\[
h_1=\operatorname{ReLU}(x-y),\qquad h_2=\operatorname{ReLU}(y-x).
\]
Sa matrice de poids est $\begin{pmatrix}1&-1\\-1&1\end{pmatrix}$ et ses deux biais sont nuls. Le neurone de sortie calcule $S=h_1+h_2$ : ses poids valent $(1,1)$ et son biais vaut $0$.}

\item \question{Calculer $S(1,1)$, $S(4,1)$ et $S(0,2)$. Une alerte se déclenche lorsque $S>2$. Pour quelles mesures se déclenche-t-elle ?}
\reponse{On trouve $S(1,1)=0$, $S(4,1)=3$ et $S(0,2)=2$. La condition étant strictement supérieure à $2$, seule la paire $(4,1)$ déclenche l'alerte.}
\end{enumerate}
}
